\chapter{Gauge--ghost и KK completion-gate}

Предыдущий heat-kernel гейт зафиксировал конечный zero-mode вклад
\begin{equation}
 B_0^{\mathrm{fin}}=\frac5{8\pi^2}
 =\frac{40}{64\pi^2}.
\end{equation}
Теперь необходимо проверить, определяет ли уже построенный parent action
оставшиеся gauge--ghost и nonzero KK слагаемые без новых нормировок.

\section{Gauge--ghost ledger}

Пусть после выбора физического gauge quotient возникает массивная
абелева векторная мода с
\begin{equation}
 c_A=\frac{m_A^2}{\chi^2}.
\end{equation}
В gauge-independent supertrace на стационарном фоне vector, Goldstone и
Faddeev--Popov блоки объединяются в три физические поляризации. Поэтому
их zero-mode вклад имеет вид
\begin{equation}
 \Delta N_{\mathrm{gauge+ghost}}=3c_A^2\geq0.
\end{equation}
Если dilaton--radion Hessian содержит дополнительную устойчивую вещественную
моду с \(c_\sigma=m_\sigma^2/\chi^2\), она добавляет \(c_\sigma^2\).

Однако алгебра \(A_F=\mathbb C\oplus\mathbb C\), first-order квадрат и
flattening-потенциал ещё не задают одновременно:
\begin{enumerate}
 \item физический quotient группы \(U(A_F)\) после действия \(J\);
 \item наличие или отсутствие динамической относительной \(U(1)\)-моды;
 \item канонически нормированный gauge coupling;
 \item BV/BRST gauge-fixing fermion и правило нулевых мод.
\end{enumerate}
В частности, текущим данным совместимы как spectator-ветвь \(c_A=0\), так
и Higgsed relative-\(U(1)\) ветвь \(c_A>0\). Во второй ветви величина
\(c_A\) зависит от ещё не выведенной канонической gauge-нормировки.

\begin{proposition}[Условная положительность gauge completion]
После задания unitary massive-vector completion его gauge--ghost вклад не
может изменить знак finite seed в отрицательную сторону. Но его величина
не является предсказанием текущего parent action, поскольку \(c_A\) не
зафиксирован.
\end{proposition}

\section{Общий KK supertrace}

Для вида поля \(s\) обозначим через \(w_s\) supertrace-вес, через
\(d_{s,n}\) кратность, а через \(\lambda_{s,n}\) безразмерный eigenvalue
соответствующего оператора на компактном носителе. Если
\begin{equation}
 \frac{M_{s,n}^2}{\chi^2}=\lambda_{s,n}+c_s,
\end{equation}
то формальный tower ledger равен
\begin{equation}
 N_{\mathrm{KK}}^{\mathrm{formal}}
 =\sum_s w_s\sum_{n\ne0}d_{s,n}
   (\lambda_{s,n}+c_s)^2.
\end{equation}
Эта сумма расходится. После общей дзета/heat-kernel процедуры требуется
конечная величина
\begin{equation}
 \Delta N_{\mathrm{KK}}
 =\operatorname{FP}_{\mathcal R}
 \sum_s w_s\sum_{n\ne0}d_{s,n}
   (\lambda_{s,n}+c_s)^2,
\end{equation}
где схема \(\mathcal R\) должна быть частью parent action, а не выбираться
после вычисления знака.

Текущая конструкция не задаёт эту сумму. Причина структурна: факторный
модуль \(\mathcal H_{\mathrm{fac}}\) был определён спектральной проекцией на
нулевые spectator-моды и не является полным локальным пространством полей
на \(K\). Поэтому scalar, spinor, one-form и ghost башни тома II нельзя
автоматически прикрепить к трём quotient-скалярам и двум Dirac-парам
версии III. Для этого нужен явный lift finite-бимодуля в полный product
Hilbert space и его field-specific Hessian complex.

\begin{nogocriterion}[Запрет наследования KK-башен]
Назначение уже известных scalar, coexact или spinor спектров тома II
полям версии III без вывода из одного fluctuated product Dirac/BV
оператора считается новой модельной гипотезой. Полученный таким способом
знак \(B\) не закрывает parent-action гейт.
\end{nogocriterion}

\section{Точный статус полного коэффициента}

На текущем уровне единственная корректная формула имеет вид
\begin{equation}
 \boxed{
 B_{\mathrm{full}}
 =\frac{40+3c_A^2+c_\sigma^2+\Delta N_{\mathrm{KK}}}
 {64\pi^2}.}
\end{equation}
Число \(40\) выведено. Остальные три величины ещё не определены одной
функциональной мерой. Поскольку \(\Delta N_{\mathrm{KK}}\) не имеет
доказанного знака, из текущих аксиом не следует ни
\(B_{\mathrm{full}}>0\), ни \(B_{\mathrm{full}}\leq0\).

\begin{theorem}[Gauge--KK specification gap]
Версия III уже содержит положительный и нормально фиксированный finite
Coleman--Weinberg seed, но ещё не определяет полный quantum scale
functional. Gauge--ghost completion является неотрицательной после выбора
устойчивой unitary ветви, тогда как величина и знак regulated KK correction
остаются неопределёнными из-за отсутствия полного fluctuated product
complex.
\end{theorem}

Это не возвращает программу к состоянию тома II. Во втором томе не был
построен общий finite parent block; теперь такой блок, его vacuum и его
finite Hessian уже фиксированы. Открытый разрыв локализован точнее: требуется
не новая численная формула, а единственный lift этого блока в gauge-fixed
локальную теорию на \(K\).

\begin{protocol}[Следующий конструктивный гейт]
Необходимо задать полный fluctuated product complex
\begin{equation}
\bigl(\mathcal A,\mathcal H,\mathcal D,J,\gamma;
  \mathcal C_{\mathrm{BV}})
\end{equation}
и из него вывести:
\begin{enumerate}
 \item physical gauge quotient и канонический coupling;
 \item scalar, vector, ghost и fermion Hessian-операторы;
 \item spin/flat-bundle ветви и правила удаления нулевых мод;
 \item одну общую конечную часть \(\Delta N_{\mathrm{KK}}\).
\end{enumerate}
Только после этого разрешён окончательный тест
\(B_{\mathrm{full}}>0\).
\end{protocol}

Машинный параметрический аудит сохранён в
\path{s2t_v3_gauge_ghost_kk_completion_results.json}.